\chapter*{From implementation to a release claim}
\addcontentsline{toc}{chapter}{Introduction: From implementation to a release claim}

Release material answers a different set of questions from a design proposal. What changed for a user of a particular version? Which contracts or defaults moved? What must an existing application do to upgrade? Which checks and artifacts support publishing that version?

The chapters in this part retain release packets, change narratives, migration guides and release-mechanics notes from the source collection. Read each chapter with its version and date in view. A guide written for an earlier upgrade describes that transition; it does not automatically describe the steps required by a later release.

\section*{Use the records together}

A release narrative describes the intended user-facing change. Migration guidance explains the work that change asks of an adopter. A release packet can link the implementation, review and validation evidence behind the claim. The design records in the preceding part provide the earlier reasoning and boundaries.

When a statement matters to an upgrade, follow it through those records and verify the package version being used. Distinguish a published package from a prerelease, a completed validation step from a planned gate, and a recorded limitation from a promised follow-up. An absence of a warning in a historical release note is not a fresh compatibility check.

\section*{Scope and live operations}

Roadmap counts, milestone assignments and tracker status can change after a document is written. Their historical wording remains useful for understanding a release decision, while the tracker remains the place to establish current scope. The same distinction applies to workflow and publishing instructions preserved in the book: current repository agreements and authorized release tooling govern a new publication.

This part is most useful as a bridge. Read from the behaviour an application relies on, through the version that changed it, to the implementation and evidence that made the change acceptable.
